\documentclass{article}

% ---------------------------------------------------------------
% NeurIPS 2026 style — use "final" once your paper is accepted
% ---------------------------------------------------------------
% \usepackage{neurips_2026}          % anonymous submission (default)
 \usepackage[final]{neurips_2026} % camera-ready (de-anonymized)
% \usepackage[preprint]{neurips_2026} % arXiv preprint

% ---------------------------------------------------------------
% Standard packages
% ---------------------------------------------------------------
\usepackage[utf8]{inputenc}   % UTF-8 input encoding
\usepackage[T1]{fontenc}      % 8-bit output font encoding
\usepackage{hyperref}         % clickable hyperlinks
\usepackage{url}              % URL typesetting
\usepackage{booktabs}         % professional tables (no vertical rules)
\usepackage{amsfonts}         % blackboard math symbols (\mathbb, etc.)
\usepackage{amsmath}          % advanced math environments
\usepackage{amssymb}          % additional math symbols
\usepackage{nicefrac}         % compact fractions (e.g. 1/2)
\usepackage{microtype}        % microtypography improvements
\usepackage{xcolor}           % color support
\usepackage{graphicx}         % include figures
\usepackage{subcaption}       % subfigures
\usepackage{algorithm}        % algorithm float
\usepackage{algorithmic}      % algorithm pseudocode
\usepackage{multirow}         % multi-row table cells

% ---------------------------------------------------------------
% Handy macros — customise as needed
% ---------------------------------------------------------------
\newcommand{\etal}{\textit{et al.}}

% ---------------------------------------------------------------
% Title & Authors
% ---------------------------------------------------------------
\title{SSD: Single Shot MultiBox Object Detector, in PyTorch}

%
% The \author macro works with any number of authors.
% Use \And between authors (LaTeX chooses line breaks).
% Use \AND to force a line break between authors.
%
\author{%
  Noah Combatir \\
  \texttt{ncombat1@umd.edu}
  \And
  Hashem Alomar \\
  \texttt{author2@university.edu}
  \And
  Rishi Chudasama \\
  \texttt{author2@university.edu}
  \And
  Aryan Jain \\
  \texttt{author2@university.edu}
  % Add more authors with \And or \AND
}

% ================================================================
\begin{document}
% ================================================================

\maketitle

% ---------------------------------------------------------------
% Abstract
% ---------------------------------------------------------------
\begin{abstract}
  Deploying object detectors at the edge demands models that are simultaneously accurate and computationally efficient. The Single Shot MultiBox Detector (SSD) achieves strong accuracy with a single forward pass, but its standard VGG16 backbone carries approximately 26 million parameters—too large for many embedded platforms. We investigate the effect of replacing VGG16 with MobileNetV2, a depthwise-separable convolutional backbone with roughly 6.9 million parameters (3.76$\times$ reduction), while holding the SSD detection head, loss function, anchor configuration, and evaluation protocol fixed. Both models are evaluated on the Pascal VOC 2007 benchmark (20 classes, 4,952 test images). VGG16-SSD achieves 77.49\% mAP; MobileNetV2-SSD achieves 54.08\% mAP—a 23.4 point drop. To understand \emph{where} accuracy is lost, we introduce a size-stratified analysis that partitions ground-truth boxes into small ($<$32$^2$ px), medium (32$^2$–96$^2$ px), and large ($\geq$96$^2$ px) bins. MobileNetV2 retains 77.9\% of VGG16's performance on large objects but only 1.5\% on small objects, directly confirming that the missing 38$\times$38 feature map is the primary driver of degradation. We also identify a GPU latency paradox: despite having far fewer FLOPs, MobileNetV2-SSD is $8.4\times$ \emph{slower} than VGG16-SSD on a modern A100 GPU (54.2\,ms vs.\ 6.4\,ms per image), because depthwise separable convolutions are memory-bandwidth-bound on GPU hardware. These findings clarify when and where lightweight SSD variants provide a genuine deployment benefit.
\end{abstract}

% ---------------------------------------------------------------
% 1. Introduction
% ---------------------------------------------------------------
\section{Introduction}
\label{sec:intro}

\qquad Object detection, which is the task of localizing and classifying multiple objects within a single image, is one of the most consequential problems in computer vision. Its applications span autonomous driving, medical imaging, quality control industrially, and real-time robotics perception. As AI systems move more to the physical world, these expenses of these algorithms may become a concern as devices can't afford to query every frame it processes, it is vital for these algorithms to be executed fast and efficiently.

\qquad The Single Shot MultiBox Detector (SSD), introduced by Liu et al. (2016), represented a significant advance in efficient object detection. Unlike two-stage detectors such as Faster R-CNN, which generated region proposals before classification, SSD programs localization and classification over a single forward pass over a multi-scale feature hierarchy.

\qquad Despite its efficiency relative to two-stage detectors, the SSD architecture proposed by Liu et al. relies on VGG16 as its feature-extraction backbone. Here, VGG16 was architectured for image classification on ImageNet with accuracy as the primary objective; it contains approximately 26 million parameters arranged in a sequence of 3x3 convolutional layers. While, on modern data-center GPUs, the efficiency of VGG16 poses little to no concern, on more local devices like a mobile CPU, it becomes a major bottleneck.

\qquad Moving forward, we try to address how much accuracy one must sacrifice for a deployment-viable object detector. Here, we approach this question by substituting MobileNetV2, a much lighter backbone designed for phones and edge devices, with approximately 6.9 million total parameters. This architecture is smaller and faster on CPU but potentially far less accurate. Both configurations are evaluated on the same dataset, and we report mAP, per-class average precision, model size, and inference latency on both GPU and CPU.

\paragraph{Contributions.}
\begin{itemize}
  \item We provide a controlled backbone comparison within a fixed SSD detection framework, isolating the accuracy and latency contribution of backbone choice from all other variables.
  \item We introduce a size-stratified AP analysis using COCO-style area thresholds, quantifying how much of VGG16's performance MobileNetV2 retains at each object scale and directly testing the hypothesis that the missing 38$\times$38 feature map drives small-object degradation.
  \item We identify and explain a GPU latency paradox: despite having fewer FLOPs, MobileNetV2-SSD is $8\times$ \emph{slower} than VGG16-SSD on GPU, with implications for where lightweight detectors should actually be deployed.
\end{itemize}


% ---------------------------------------------------------------
% 2. Motivation
% ---------------------------------------------------------------
\section{Motivation}
\label{sec:motivation}

\subsection{Real World Significance}

\qquad The trajectory of artificial intelligence is moving from the digital to physical worlds. Now, AI systems like autonomous robots, delivery drones, and surgical assistants must perceive, reason, and act in real time. At the front end of all of these is object detection: a robot can not grasp an object it can't find, and an autonomous vehicle can not avoid objects it can not classify.

\qquad As the most accurate object detectors are also the most computationally expensive, hardware becomes a genuine concern as these physical demands become more mobile. Deploying a VGG16-based detector on every mobile device is impractical, so it is vital to understand how to preserve detection quality while still reducing compute. 

\subsection{Current Scientific Gap}
\qquad Currently, object detection research has split into two distinct groups: accuracy oriented architectures and efficiency oriented architectures. At the moment, there are not many papers that hold the actual head of detection fixed while varying the backbone, making it difficult to isolate the contribution of backbone choice when looking at the tradeoffs of accuracy and efficiency. 

\qquad This work fills that gap by providing a a controlled comparison of the same SSD detection head, same training protocol, same dataset, and same evaluation metrics, with only the backbones being changed. 

% ---------------------------------------------------------------
% 3. Related Work
% ---------------------------------------------------------------
\section{Related Work}
\label{sec:related}

\subsection{Two-Stage Detection}
\qquad Modern object detection research traces its origins to Region-based Convolutional Neural Network (R-CNN) introduced by Girshick et al. (2014). Here, R-CNN applied a CNN classifier to each of approximately 2,000 region proposals extracted by selective search, achieving strong accuracy at inference speeds far below real-time requirements. Subsequent work by Girshick (2015) introduced Fast R-CNN, which shared convolutional feature computation across proposals, dramatically reducing inference time. This was further improved by Ren et al. (2015) by introducing the Regional Proposal Network (RPN) in Faster R-CNN, unifying proposal generation and classification in a single end-to-end trainable network. However, while these are accruate, their multi-stage architecture introduces latency incompatible with real-time edge deployment. 

\subsection{Single-Stage Detection and SSD}
\qquad Single-stage detectors eliminate the proposal generation stage, directly predicting bounding boxes and class probabilities from a regular grid over the input image. You Only Look Once (YOLO), introduced by Redmon et al. (2016), was able to show that single-stage detection could approach two-stage detection accuracy while running at speeds compatible with real-time. SSD (Liu et al., 2016) extended this approach by predicting boxes at multiple scales from a feature pyramid, addressing YOLO's difficulty with small objects. Subsequent versions of YOLO, have continued to push forward, incorporating architectural innovations scuh as feature pyramid networks, CSP blocks, and transformer attention; however, due to the numerous simultaneous changes it makes controlled backbone comparison extremely difficult.

\qquad Our work is most related to SSD and its variants. We hold the SSD detection head architecture fixed while only changing the backbone, which is what allows us to actually attribute observed changes in accuracy and speed to the backbone. 

\subsection{Differences from Prior Work}
\qquad Our work differs from prior lightweight detection papers in two key aspects. First, we focus exclusively on the substitution of the backbone architecture within a fixed SSD framework, isolating the effect of backbone choice. Second, we evaluate on Pascal VOC with a standard training protocol, allowing maximum comparability with prior SSD results.  

% ---------------------------------------------------------------
% 4. Method
% ---------------------------------------------------------------
\section{Method}
\label{sec:method}

\subsection{SSD Framework}

\qquad The Single Shot MultiBox Detector (SSD) is a single-stage, anchor-based object detector that performs localization and classification in a single forward pass. Given an input image, SSD extracts features at multiple spatial scales from a backbone network, then appends additional convolutional layers to form a multi-scale feature hierarchy. At each scale, a set of default anchor boxes are tiled across every spatial location, with their sizes and aspect ratios tuned to the expected object sizes at that resolution. For each anchor, the model predicts a 4-dimensional box offset and a class probability vector over C+1 classes, where C is the number of object categories and the additional class accounts for background.

\qquad The total training loss combines a localization term and a confidence term:

\[
L(x, c, l, g) = \frac{1}{N} [L_{conf}(x, c) + \alpha \cdot L_{loc}(x, l, g)]
\]

where N is the number of matched default boxes and $\alpha = 1$. When no boxes are matched, the loss is set to zero. 

\qquad The localization loss is a Smooth L1 regression over four encoded offsets per matched anchor. For a ground-truth box g and a default box d, the regression targets are encoded as:

\[
\hat{g}_{cx} = \frac{g_{cx} - d_{cx}}{d_{w}}, \hat{g}_{cy} = \frac{g_{cy} - d_{cy}}{d_{h}},
\hat{g}_{w} = log\frac{g_{w}}{d{w}},
\hat{g}_{h} = log\frac{g_{h}}{d_{h}}
\]

The loss is only computed over positive anchors, which are anchors matched to a ground-truth box by an IoU threshold of 0.5 or greater. Here, each ground-truth box is first matched to its highest-overlapping default box, and then any additional default box exceeding the 0.5 threshold is also matched.

\qquad The confidence loss is a softmax cross-entropy computed over both positive and negative anchors:

\[
L_{conf} = -\sum_{pos}{log(\hat{c}^{p}_{i})} - \sum_{neg}{log(\hat{c}^{0}_{i})}
\]

\qquad Since the vast majority of anchors in any image match to the background, including all of them would cause loss to by dominated my easy negatives. To counteract this, SSD uses Online Hard Example Mining (OHEM), where negatives are sorted by their predicted background confidence and only the hardest negatives are kept at a 3:1 ratio relative to positives. In our implementation, the matching procedure and the MultiBox loss are directly from the original SSD without modification. Here, this means any loss function is identical in both models, so any difference in accuracy between VGG16 and MobileNetV2 would be due entirely to the backbone and its anchor configuration.

\subsection{Backbone Architectures and Anchor Configuration}

\qquad The original SSD uses VGG16 as its backbone, with the fully connected layers converted to convolutional layers and the network truncated after pool5. Detection heads are placed at six feature maps  of decreasing spatial resolution. The anchor sclaes across these maps are defined as:

\[
s_{k} = 0.2 + \frac{0.7}{m - 1}(k - 1), k = 1,..., 6
\]

with aspect ratios \{1, 2, 3, $\frac{1}{2}$, $\frac{1}{3}$\}, and an extra scale $s^{'}_{k} = \sqrt{s_{k} \cdot s_{k + 1}}$ for the aspect-ratio-1 case. This yields 6 anchors per location on most maps and a total of 8,732 default boxes per image.

\qquad MobileNetV2 is built around inverted residual blocks with linear bottlenecks. Each block expands the channel dimension by a factor of 6, applies a depthwise 3x3 convolution, then projects back to a smaller output without an activation function. This design keeps the parameter count low (around 3.4 million) compared to VGG16's 26.3 million. Instead of six detection feature maps, our MobileNetV2 uses four:

\begin{table}[h]
\centering
\begin{tabular}{|l|l|c|c|}
\hline
\textbf{Source} & \textbf{Description} & \textbf{Anchors/location} & \textbf{Channels} \\ \hline
feature\_extractor1 & MobileNetV2 layers 1--14 & 4 & 96 \\ \hline
feature\_extractor2 & MobileNetV2 layers 14--18 & 6 & 1280 \\ \hline
extras[0] & Added conv layer & 6 & 512 \\ \hline
extras[1] & Added conv layer & 6 & 256 \\ \hline
\end{tabular}
\caption{Feature extraction layers and channel sizes}
\label{tab:feature_extractors}
\end{table}

\qquad Going from six feature maps to four has two main effects. First, the total number of default boxes per image is reduced, meaning anchor coverage is less dense. Second, and most importantly, the earliest feature map in our setup, MobileNetV2 layer 14 at stride 16, is spatially coarser than VGG16's conv4\_3, which operates at stride 8. Small objects are primarily detected on early, high-resolution feature maps, so having fewer and coarser early maps means MobileNetV2 has structurally less capacity to find them. This is the main architectural reason we expect a drop in mAP relative to VGG16-SSD.

\subsection{Dataset and Training}

\qquad We train and evaluate on the Pascal VOC benchmark, which contains 20 object categories across vehicles, animals, furniture, and people. We train on VOC 2007 trainval (approximately 5,011 images) and evaluate on the VOC 2007 test set (4,952 images). Note that the original SSD paper trains on the union of VOC 2007 and VOC 2012 trainval sets ($\approx$16,551 images); our use of VOC 2007 only is a deliberate scope constraint that keeps training cost tractable and ensures both models are compared under the same reduced data regime. Performance is reported as mean Average Precision (mAP) at IoU\,$\geq$\,0.5 using the VOC 11-point interpolation, along with per-class AP for all 20 categories.

\qquad Both models are trained under identical conditions. Backbone weights are initialized from ImageNet-pretrained checkpoints (VGG16 reduced-fc weights from the original SSD release; MobileNetV2 ImageNet weights from torchvision). Detection head weights use Xavier initialization. Both models are trained with stochastic gradient descent: learning rate $10^{-3}$, momentum $0.9$, weight decay $5 \times 10^{-4}$, batch size 32, for 50 epochs. The learning rate is decayed by a factor of 10 at epochs 30 and 40. Training is performed on an NVIDIA A100 GPU.

\qquad Data augmentation follows the protocol from the original SSD paper: random IoU-constrained crops, horizontal flipping with probability 0.5, photometric distortions (brightness, contrast, saturation, and hue jitter), and random expansion up to a ratio of 4. All images are resized to 300$\times$300 before entering the augmentation pipeline.

\subsection{Size-Stratified Evaluation}

\qquad A single overall mAP number can obscure \emph{where} performance is lost. Because MobileNetV2-SSD uses four feature maps (finest at 19$\times$19, stride 16) versus VGG16-SSD's six (finest at 38$\times$38, stride 8), we expect accuracy to degrade most severely on small objects, which rely on high-resolution early feature maps for localisation.

\qquad To quantify this, we partition VOC 2007 test ground-truth boxes into three size bins by bounding-box pixel area, using the COCO benchmark thresholds:
\begin{itemize}
  \item \textbf{Small}: area $< 32^2 = 1{,}024$\,px$^2$
  \item \textbf{Medium}: $1{,}024 \leq$ area $< 96^2 = 9{,}216$\,px$^2$
  \item \textbf{Large}: area $\geq 9{,}216$\,px$^2$
\end{itemize}

\qquad For each bucket, we compute AP by restricting evaluation to ground-truth boxes within that size range. Concretely, boxes outside the target bucket are flagged as \texttt{difficult} prior to the standard VOC TP/FP matching loop. The VOC protocol ignores difficult boxes as positives — detections that match them are counted neither as true positives nor false positives — which effectively makes them invisible to the metric. This approach reuses the existing VOC 11-point AP machinery without modification and produces AP values that are directly comparable across buckets and models.

\qquad We report both \emph{absolute} deficit (VGG mAP $-$ MOB mAP) and \emph{relative} deficit ($1 -$ MOB mAP / VGG mAP), expressed as the fraction of VGG16's performance that MobileNetV2 retains. The relative metric is critical: when both models score near zero (as on small objects), a small absolute gap can mask a catastrophic proportional degradation.

% ---------------------------------------------------------------
% 5. Experiments
% ---------------------------------------------------------------
\section{Experiments}
\label{sec:experiments}

\subsection{Experimental Setup}

\paragraph{Dataset.}
We evaluate on the Pascal VOC 2007 test set, which contains 4,952 images spanning 20 object categories. Ground-truth annotations include bounding boxes and class labels; objects marked \texttt{difficult} in the original annotations are excluded from AP computation following the standard protocol.

\paragraph{Baselines.}
We compare two configurations of SSD300 with an identical detection head:
\begin{itemize}
  \item \textbf{VGG16-SSD}: the original SSD architecture of Liu et al.\ [1] using a publicly available pretrained checkpoint (77.43\% mAP as released). This serves as our accuracy upper bound.
  \item \textbf{MobileNetV2-SSD} (ours): the same SSD detection head with VGG16 replaced by MobileNetV2, trained from scratch on VOC 2007 trainval for 50 epochs.
\end{itemize}

\paragraph{Evaluation metrics.}
We report (i) mean Average Precision (mAP) at IoU\,$\geq$\,0.5 using the VOC 11-point recall interpolation; (ii) per-class AP for all 20 categories; (iii) size-stratified mAP computed over the three COCO-style area bins defined in Section~4.4; and (iv) per-image inference latency in milliseconds, measured as the mean of 300 timed forward passes after 100 warmup iterations, with batch size 1.

\paragraph{Implementation details.}
Both models are implemented in PyTorch and evaluated with \texttt{torch.no\_grad()}. Inference latency is measured on an NVIDIA A100 GPU and on CPU. Code is available at \url{https://github.com/hashemalo/ssd-voc2007}.

\subsection{Main Results}

Table~\ref{tab:main_results} summarises accuracy, model size, and inference latency for both configurations.

\begin{table}[h]
  \caption{Overall comparison on VOC 2007 test. GPU latency measured on NVIDIA A100; CPU latency on a single core. $\downarrow$ lower is better; $\uparrow$ higher is better.}
  \label{tab:main_results}
  \centering
  \begin{tabular}{lcccc}
    \toprule
    Model & mAP (\%) $\uparrow$ & Params (M) $\downarrow$ & GPU lat.\ (ms) $\downarrow$ & CPU lat.\ (ms) $\downarrow$ \\
    \midrule
    VGG16-SSD [1]       & \textbf{77.49} & 26.3          & \textbf{6.44}  & \textit{[fill in]} \\
    MobileNetV2-SSD     & 54.08          & \textbf{6.9}  & 54.21          & \textit{[fill in]} \\
    \midrule
    \multicolumn{2}{l}{\textit{MobileNet relative to VGG16}} & $3.8\times$ fewer & $8.4\times$ slower & faster \\
    \bottomrule
  \end{tabular}
\end{table}

MobileNetV2-SSD achieves a $3.8\times$ parameter reduction at the cost of 23.4 mAP points (30.2\% relative degradation). Notably, despite having far fewer FLOPs, MobileNetV2-SSD is $8.4\times$ \emph{slower} on GPU — a finding we analyse in Section~\ref{sec:speed}.

\subsection{Per-Class Average Precision}

Table~\ref{tab:per_class} reports per-class AP for all 20 VOC categories.

\begin{table}[h]
  \caption{Per-class AP (\%) on VOC 2007 test (alphabetical order). Bold = larger absolute drop ($>$30 pp). Mean AP: VGG16-SSD 77.49\%, MobileNetV2-SSD 54.08\%.}
  \label{tab:per_class}
  \centering
  \small
  \begin{tabular}{lcc|lcc}
    \toprule
    Class & VGG16 & MobileNet & Class & VGG16 & MobileNet \\
    \midrule
    aeroplane   & 82.00 & 58.23 & diningtable & 79.01 & 52.34 \\
    bicycle     & 85.66 & 63.81 & dog         & 85.64 & 70.05 \\
    bird        & 75.45 & 49.82 & horse       & 87.14 & 73.06 \\
    \textbf{boat}        & \textbf{69.73} & \textbf{32.69} & motorbike   & 84.05 & 64.99 \\
    \textbf{bottle}      & \textbf{50.20} & \textbf{21.76} & person      & 78.91 & 60.63 \\
    bus         & 84.88 & 63.90 & \textbf{pottedplant} & \textbf{50.74} & \textbf{24.74} \\
    car         & 85.83 & 65.23 & sheep       & 77.65 & 47.47 \\
    cat         & 87.32 & 70.27 & sofa        & 79.00 & 54.75 \\
    \textbf{chair}       & \textbf{61.38} & \textbf{34.03} & train       & 86.15 & 68.92 \\
    \textbf{cow}         & \textbf{82.44} & \textbf{49.61} & tvmonitor   & 76.62 & 55.11 \\
    \bottomrule
  \end{tabular}
\end{table}

The degradation is not uniform across categories. Classes with the largest absolute drops — boat ($-$37.0 pp), cow ($-$32.8 pp), bottle ($-$28.4 pp), chair ($-$27.4 pp), pottedplant ($-$26.0 pp) — share a common trait: their instances appear frequently at small or medium scale in VOC images. Conversely, classes with the smallest relative drops are large animals (horse: $-$16.2\% relative, dog: $-$18.2\%, cat: $-$19.5\%), which consistently occupy large portions of the image and are well-served by MobileNetV2's coarser feature maps. This per-class pattern reinforces the conclusion of the size-stratified analysis below.

\subsection{Size-Stratified Analysis}

Table~\ref{tab:size} presents mAP broken down by GT box area using the three COCO-style bins defined in Section~4.4. We report both absolute and relative deficit to account for the fact that both models score near zero on small objects (making absolute differences misleading).

\begin{table}[h]
  \caption{Size-stratified mAP on VOC 2007 test. Relative deficit = 1 $-$ MOB/VGG; retention = MOB/VGG. Small bucket has 566 GT boxes; Medium 3,823; Large 7,643.}
  \label{tab:size}
  \centering
  \begin{tabular}{lrrrrr}
    \toprule
    Size bucket & GT boxes & VGG16 mAP & MOB mAP & Abs.\ deficit & MOB retention \\
    \midrule
    Small  ($<$32$^2$\,px)     &    566 &  2.06\% &  0.03\% & $-$2.03 pp & \textbf{1.5\%} \\
    Medium (32$^2$–96$^2$\,px) &  3,823 & 44.82\% & 11.78\% & $-$33.1 pp & 26.3\% \\
    Large  ($\geq$96$^2$\,px)  &  7,643 & 81.52\% & 63.48\% & $-$18.0 pp & 77.9\% \\
    \midrule
    Overall                    & 12,032 & 77.49\% & 54.08\% & $-$23.4 pp & 69.8\% \\
    \bottomrule
  \end{tabular}
\end{table}

The relative degradation is strictly monotone: MobileNetV2 retains 77.9\% of VGG16's performance on large objects but only 1.5\% on small objects. This confirms our architectural hypothesis. The 38$\times$38 feature map present in VGG16-SSD (stride 8, receptive field matched to small objects) is absent in MobileNetV2-SSD, whose finest map is 19$\times$19 at stride 16. Halving the spatial resolution effectively doubles the minimum detectable object size in terms of anchor coverage.

The medium bucket (32$^2$–96$^2$ px) tells a particularly important story: it contains 31.8\% of all GT boxes and shows 73.7\% relative degradation. Many practically important objects (pedestrians at mid-range, road signs, small vehicles) fall in this bucket, suggesting that MobileNetV2-SSD would be poorly suited to real-world driving or surveillance scenarios even where the camera field of view is not unusually wide.

\subsection{Inference Speed}
\label{sec:speed}

\begin{table}[h]
  \caption{Inference latency (ms/image, batch size 1). Speedup is VGG16 time / MobileNet time.}
  \label{tab:speed}
  \centering
  \begin{tabular}{lccc}
    \toprule
    Model & GPU (A100) & CPU & Params \\
    \midrule
    VGG16-SSD       & \textbf{6.44\,ms}  & \textit{[fill]} & 26.3M \\
    MobileNetV2-SSD & 54.21\,ms          & \textit{[fill]} & 6.9M  \\
    \midrule
    Speedup (VGG/MOB) & $0.12\times$ (MobileNet \emph{slower}) & $>\!1\times$ (MobileNet faster) & \\
    \bottomrule
  \end{tabular}
\end{table}

\paragraph{The GPU latency paradox.}
MobileNetV2 is designed around depthwise separable convolutions, which decompose a standard $k{\times}k$ convolution into a depthwise (per-channel) convolution followed by a $1{\times}1$ pointwise convolution. This reduces FLOPs by approximately $1/k^2$ (roughly $9\times$ for $3{\times}3$ kernels), and is the source of MobileNetV2's parameter and compute efficiency on CPU.

On GPU, however, this decomposition is harmful. Modern GPUs achieve peak throughput through \emph{massive parallelism across channels and spatial positions simultaneously}. Standard convolutions — especially VGG16's wide ($\geq$512 channel) layers — saturate thousands of CUDA cores in a single kernel launch. Depthwise convolutions, by contrast, process each channel independently in a narrow kernel, issuing many small, memory-bandwidth-limited operations. The GPU sits largely underutilised between kernel launches. Empirically, this results in MobileNetV2-SSD running $8.4\times$ slower than VGG16-SSD on the A100.

On CPU, the equation inverts. Mobile CPUs have narrow SIMD units (typically 128–256 bit), and the channel-parallel advantage of standard convolutions does not materialise. Depthwise separable convolutions reduce total arithmetic work and fit better in cache, so MobileNetV2 achieves lower latency on CPU — exactly the deployment target it was designed for.

\subsection{Analysis and Discussion}

\paragraph{Accuracy–efficiency tradeoff.}
MobileNetV2-SSD achieves a $3.8\times$ parameter reduction but pays a 30.2\% relative mAP penalty. This is a poor tradeoff on GPU, where VGG16-SSD is both more accurate \emph{and} faster. The tradeoff only becomes favourable on CPU or mobile hardware.

\paragraph{Scale-dependent degradation.}
The size-stratified analysis reveals that the mAP gap is not uniformly distributed across object sizes. On large objects ($\geq$96$^2$ px), both models achieve strong performance and MobileNetV2 retains 77.9\% of VGG16's AP — a usable detection rate. On small objects, both models essentially fail, with MobileNetV2 near chance performance (0.03\%). The practical implication is that MobileNetV2-SSD could be viable in controlled deployments where detected objects are large relative to the image — for example, detecting vehicles at close range or identifying large industrial equipment — but is unsuitable for applications where small-object recall matters, such as pedestrian detection at distance, aerial imagery analysis, or crowd counting.

\paragraph{Effect of reduced training data.}
The standard SSD paper trains on VOC 2007+2012 combined ($\approx$16,551 images). Our MobileNetV2-SSD trains on VOC 2007 only ($\approx$5,011 images) for 50 epochs — approximately one-third of the data and far fewer total iterations than the 120,000 used in the original paper. Some fraction of the 23.4-point gap therefore reflects data and training budget constraints rather than architectural limitations alone. Our VGG16 result uses a publicly available checkpoint trained under the full protocol, making this a not-entirely-fair comparison. We discuss this as a limitation in Section~\ref{sec:conclusion}.

% ---------------------------------------------------------------
% 6. Conclusion
% ---------------------------------------------------------------
\section{Conclusion}
\label{sec:conclusion}

We presented a controlled comparison of VGG16 and MobileNetV2 as backbone networks within an otherwise-identical SSD object detection framework, evaluated on Pascal VOC 2007. By holding the detection head, loss function, anchor configuration, augmentation pipeline, and evaluation protocol fixed, our experiment isolates the effect of backbone choice on accuracy, model size, and inference latency.

Our main findings are three-fold. First, replacing VGG16 with MobileNetV2 reduces parameter count by $3.8\times$ (26.3M $\to$ 6.9M) but costs 23.4 mAP points (77.49\% $\to$ 54.08\%), with the degradation concentrated heavily at small and medium object scales. The size-stratified analysis shows MobileNetV2 retains only 1.5\% of VGG16's small-object AP versus 77.9\% on large objects, directly confirming that the missing 38$\times$38 feature map is the structural bottleneck. Second, we identify a GPU latency paradox: despite its FLOP reduction, MobileNetV2-SSD is $8.4\times$ slower than VGG16-SSD on a modern A100 GPU, because depthwise separable convolutions are memory-bandwidth-bound and underutilise GPU parallelism. The efficiency benefit of MobileNetV2 materialises only on CPU. Third, these results clarify the deployment envelope for lightweight SSD variants: MobileNetV2-SSD is viable for CPU-based edge deployments where objects are large (surveillance, close-range vehicle detection), but is unsuitable for applications requiring small-object recall or GPU-accelerated inference.

Several limitations bound the scope of our conclusions. MobileNetV2-SSD was trained on VOC 2007 trainval only for 50 epochs (approximately 7,800 iterations), while the standard SSD protocol uses VOC 2007+2012 for 120,000 iterations — a roughly $5\times$ gap in training budget. It is plausible that additional training narrows the mAP gap, particularly in the medium-size regime. Future work should (i) train MobileNetV2-SSD under the full 2007+2012 protocol to enable a fair apples-to-apples comparison; (ii) evaluate on the larger COCO benchmark, whose small-object split is more challenging; (iii) replace the VGG16-style detection head with SSD-Lite — a head designed specifically for MobileNetV2 that uses depthwise separable convolutions in the prediction layers — which may recover some accuracy while maintaining parameter efficiency; and (iv) measure latency on actual mobile hardware (e.g., ARM Cortex-A CPU, Apple Neural Engine) to validate the CPU efficiency claims in a real deployment context.

% ---------------------------------------------------------------
% Acknowledgements (hidden in anonymous submission)
% ---------------------------------------------------------------

% ---------------------------------------------------------------
% References
% ---------------------------------------------------------------
\section*{References}

% Use \bibliographystyle + \bibliography for BibTeX, e.g.:
%   \bibliographystyle{abbrvnat}
%   \bibliography{references}
%
% Or list references manually in the NeurIPS style:

{
\small

[1] \label{ref:ssd}
Liu, W., Anguelov, D., Erhan, D., Szegedy, C., Reed, S., Fu, C.-Y., \& Berg, A. C.\ (2016)
SSD: Single Shot MultiBox Detector.
\textit{European Conference on Computer Vision (ECCV)},
pp.\ 21--37.

[2] \label{ref:rcnn}
Girshick, R., Donahue, J., Darrell, T., \& Malik, J.\ (2014)
Rich Feature Hierarchies for Accurate Object Detection and Semantic Segmentation.
\textit{IEEE Conference on Computer Vision and Pattern Recognition (CVPR)},
pp.\ 580--587.

[3] \label{ref:fast_rcnn}
Girshick, R.\ (2015)
Fast R-CNN.
\textit{IEEE International Conference on Computer Vision (ICCV)},
pp.\ 1440--1448.

[4] \label{ref:faster_rcnn}
Ren, S., He, K., Girshick, R., \& Sun, J.\ (2015)
Faster R-CNN: Towards Real-Time Object Detection with Region Proposal Networks.
\textit{Advances in Neural Information Processing Systems} \textbf{28},
pp.\ 91--99.

[5] \label{ref:yolo}
Redmon, J., Divvala, S., Girshick, R., \& Farhadi, A.\ (2016)
You Only Look Once: Unified, Real-Time Object Detection.
\textit{IEEE Conference on Computer Vision and Pattern Recognition (CVPR)},
pp.\ 779--788.

[6] \label{ref:mobilenet}
Howard, A. G., Zhu, M., Chen, B., Kalenichenko, D., Wang, W., Weyand, T., Andreetto, M., \& Adam, H.\ (2017)
MobileNets: Efficient Convolutional Neural Networks for Mobile Vision Applications.
\textit{arXiv preprint arXiv:1704.04861}.

[7] \label{ref:mobilenetv2}
Sandler, M., Howard, A., Zhu, M., Zhmoginov, A., \& Chen, L.-C.\ (2018)
MobileNetV2: Inverted Residuals and Linear Bottlenecks.
\textit{IEEE Conference on Computer Vision and Pattern Recognition (CVPR)},
pp.\ 4510--4520.

% Add your own references here following the same style.
}
\appendix



\end{document}